\documentclass[11pt]{article}
\usepackage[utf8]{inputenc}
\usepackage[T1]{fontenc}
\usepackage[a4paper,margin=1in]{geometry}
\usepackage{lmodern}
\usepackage{microtype}
\usepackage{amsmath,amssymb,mathtools}
\usepackage{mathrsfs}
\usepackage{tikz-cd}
\usetikzlibrary{arrows}
\usepackage{enumitem}
\usepackage{longtable,booktabs,array}
\usepackage[hidelinks]{hyperref}
\setlength{\parindent}{0pt}
\setlength{\parskip}{0.55em}
\emergencystretch=3em
\hbadness=10000
\hfuzz=20pt
\usepackage[normalem]{ulem}
\newcommand{\K}{K}
\newcommand{\Gr}{\operatorname{Gr}}
\newcommand{\Fil}{\operatorname{Fil}}
\newcommand{\ch}{\operatorname{ch}}
\newcommand{\Todd}{\operatorname{Todd}}
\newcommand{\cl}{\operatorname{cl}}
\newcommand{\Parf}{\operatorname{Parf}}
\newcommand{\Tor}{\operatorname{Tor}}
\newcommand{\Spec}{\operatorname{Spec}}
\newcommand{\RRR}{\mathrm{RRR}}
\providecommand{\codim}{\operatorname{codim}}
\providecommand{\Inv}{\operatorname{Inv}}
\providecommand{\Pic}{\operatorname{Pic}}
\providecommand{\Z}{\mathbb Z}
\providecommand{\Q}{\mathbb Q}
\providecommand{\Ahat}{\widehat A}
\providecommand{\Ahp}{\widehat A^{+}}
\providecommand{\That}{\widehat T}
\providecommand{\rad}{\operatorname{rad}}
\providecommand{\Gl}{\operatorname{Gl}}
\providecommand{\GL}{\operatorname{Gl}}
\providecommand{\ch}{\operatorname{ch}}
\providecommand{\Todd}{\operatorname{Todd}}
\providecommand{\Gr}{\operatorname{Gr}}
\providecommand{\red}{\text{réd}}

% Core macro additions
\providecommand{\Atil}{\widetilde A}
\providecommand{\Ctil}{\widetilde C}
\providecommand{\K}{\mathrm K}
\providecommand{\cC}{\mathcal C}
\providecommand{\cF}{\mathscr F}
\providecommand{\eps}{\varepsilon}
\providecommand{\ellc}{\ell}
\providecommand{\Tor}{\operatorname{Tor}}
\providecommand{\Coker}{\operatorname{Coker}}
\providecommand{\Ker}{\operatorname{Ker}}
\providecommand{\Supp}{\operatorname{Supp}}
\providecommand{\codim}{\operatorname{codim}}
\providecommand{\Gr}{\operatorname{G}}
\providecommand{\rank}{\operatorname{rank}}
\providecommand{\rang}{\operatorname{rang}}
\providecommand{\cO}{\mathcal O}
\providecommand{\sgacircled}[1]{\tikz[baseline=(char.base)]\node[shape=circle,draw,inner sep=.35pt](char){\mbox{\scriptsize #1}};}

\hypersetup{hypertexnames=false}
\newcommand{\sgastandalonefootnotetext}[2]{%
  \begingroup
  \NoHyper
  \footnotetext{#2}%
  \endNoHyper
  \endgroup
}
\providecommand{\R}{\mathbb R}
\providecommand{\C}{\mathbb C}
\providecommand{\N}{\mathbb N}
\providecommand{\Zell}{\mathbb Z_\ell}
\providecommand{\Qell}{\mathbb Q_\ell}
\providecommand{\C}{\mathbb C}
\providecommand{\Ch}{\operatorname{Ch}}
\providecommand{\Sym}{\operatorname{Sym}}
\providecommand{\Td}{\operatorname{Td}}
\begin{document}
\providecommand{\cS}{\mathcal S}
\providecommand{\cD}{\mathcal D}
\providecommand{\cE}{\mathscr E}
\providecommand{\cG}{\mathscr G}
\providecommand{\cL}{\mathscr L}
\providecommand{\cM}{\mathscr M}
\providecommand{\cP}{\mathscr P}
\providecommand{\cK}{\mathcal K}
\providecommand{\Db}{\mathrm D}
\providecommand{\Kb}{\mathrm K}
\providecommand{\Trc}{\operatorname{Tr}}
\providecommand{\ob}{\operatorname{ob}}
\providecommand{\Coh}{\operatorname{Coh}}
\providecommand{\Cone}{\operatorname{Cone}}
\providecommand{\im}{\operatorname{Im}}
\providecommand{\Id}{\operatorname{Id}}
\providecommand{\RR}{\mathrm{R.R.}}
\providecommand{\RRR}{\mathrm{R.R.R.}}
\providecommand{\GrR}{\mathrm{G}}
\providecommand{\pr}{\operatorname{pr}}
\providecommand{\cC}{\mathcal C}
\providecommand{\cD}{\mathcal D}
\providecommand{\cK}{\mathcal K}
\providecommand{\cO}{\mathcal O}
\providecommand{\cS}{\mathcal S}
\providecommand{\Cpx}{\mathrm C}
\providecommand{\Kpx}{\mathrm K}
\providecommand{\Dpx}{\mathrm D}
\providecommand{\D}{\mathrm D}
\providecommand{\Ho}{\mathrm H}
\providecommand{\RHom}{\mathrm R\!\operatorname{Hom}}
\providecommand{\Ob}{\operatorname{ob}}
\providecommand{\Coker}{\operatorname{Coker}}
\providecommand{\Ker}{\operatorname{Ker}}
\renewcommand{\Im}{\operatorname{Im}}
\providecommand{\Hom}{\operatorname{Hom}}
\providecommand{\Cone}{\operatorname{Cone}}
\providecommand{\id}{\operatorname{Id}}
\providecommand{\qis}{\xrightarrow{\sim}}
\providecommand{\RHom}{\mathrm R\!\Hom}
\providecommand{\Ext}{\operatorname{Ext}}
\providecommand{\coh}{\operatorname{coh}}
\providecommand{\PreCoh}{\operatorname{precoh}}
\providecommand{\Tr}{\operatorname{Tr}}
\providecommand{\Qis}{\operatorname{Qis}}
\providecommand{\Rad}{\operatorname{rad}}
\providecommand{\Spec}{\operatorname{Spec}}
\providecommand{\End}{\operatorname{End}}
\providecommand{\Id}{\operatorname{Id}}
\providecommand{\varinjlim}{\mathop{\mathrm{lim}}\limits_{\longrightarrow}}
\providecommand{\varprojlim}{\mathop{\mathrm{lim}}\limits_{\longleftarrow}}
\providecommand{\qcoh}{\operatorname{qcoh}}
\providecommand{\amp}{\operatorname{amp}}
\providecommand{\Cpx}{\mathrm C}
\providecommand{\Kpx}{\mathrm K}
\providecommand{\Dpx}{\mathrm D}
\providecommand{\Ho}{\mathrm H}
\providecommand{\RHom}{\mathrm R\!\operatorname{Hom}}
\renewcommand{\Ob}{\operatorname{ob}}
\providecommand{\qis}{\xrightarrow{\sim}}
\providecommand{\RHom}{\mathrm R\!\Hom}
\providecommand{\Ext}{\operatorname{Ext}}
\providecommand{\Hom}{\operatorname{Hom}}
\providecommand{\Tr}{\operatorname{Tr}}
\providecommand{\Add}{\operatorname{Add}}
\providecommand{\Cart}{\underline{\operatorname{Cart}}}
\providecommand{\perfamp}{\operatorname{parf.amp}}
\providecommand{\perfdim}{\operatorname{parf.dim}}
\providecommand{\toramp}{\operatorname{tor.amp}}
\providecommand{\torf}{\operatorname{torf}}
\providecommand{\injf}{\operatorname{injf}}
\providecommand{\rg}{\operatorname{rg}}
\providecommand{\is}{\mathrm{is}}
\providecommand{\Mod}{\operatorname{Mod}}
\providecommand{\Cpx}{\mathrm C}
\providecommand{\Kpx}{\mathrm K}
\providecommand{\Dpx}{\mathrm D}
\providecommand{\Ho}{\mathrm H}
\providecommand{\RHom}{\mathrm R\!\operatorname{Hom}}
\renewcommand{\Ob}{\operatorname{ob}}
\providecommand{\Hom}{\operatorname{Hom}}
\providecommand{\RHom}{\mathrm R\!\operatorname{Hom}}
\providecommand{\Ext}{\operatorname{Ext}}
\providecommand{\Tor}{\operatorname{Tor}}
\providecommand{\Tr}{\operatorname{Tr}}
\providecommand{\Add}{\operatorname{Add}}
\providecommand{\Cart}{\underline{\operatorname{Cart}}}
\providecommand{\Parf}{\operatorname{Parf}}
\providecommand{\coh}{\operatorname{coh}}
\providecommand{\amp}{\operatorname{amp}}
\providecommand{\cl}{\operatorname{cl}}
\providecommand{\rg}{\operatorname{rg}}
\providecommand{\Id}{\operatorname{Id}}
\providecommand{\Ker}{\operatorname{Ker}}
\renewcommand{\Im}{\operatorname{Im}}
\providecommand{\cH}{\mathcal H}
\providecommand{\Zs}{\mathbb Z_S}
\providecommand{\A}{\mathcal A}
\providecommand{\cC}{\mathcal C}
\providecommand{\cCzero}{\mathcal C_0}
\providecommand{\otL}{\mathop{\otimes}^{\mathrm L}}
\providecommand{\tamp}{\operatorname{tor.amp}}
\providecommand{\pamp}{\operatorname{parf.amp}}
\providecommand{\pdim}{\operatorname{parf.dim}}
\providecommand{\bon}{\mathrm b}
\providecommand{\parf}{\operatorname{parf}}
\providecommand{\parff}{\operatorname{parf}}
\providecommand{\ac}{\mathrm{ac}}
\providecommand{\End}{\operatorname{End}}
\providecommand{\tr}{\operatorname{tr}}
%%%%%%%%%%  scan idx 169  |  folio 156  %%%%%%%%%%
On
peut interpréter \(u\) (resp. \(v\)) comme un homomorphisme, dans \(\Dpx(U)\), de \(A|U\) dans \(\mathbb{R}\underline{\Hom}(E,F)|U\) (resp. \(\mathbb{R}\underline{\Hom}(F,E)|U\)), et former
\[
  u\otL v:A|U\longrightarrow(\mathbb{R}\underline{\Hom}(E,F)\otL\mathbb{R}\underline{\Hom}(F,E))|U.
\]
L'application \((u,v)\mapsto\langle u,v\rangle=\gamma\circ(u\otL v)\) définit une application \(H^0(U,A)\)-linéaire,
\[
  \Hom_{\Dpx(U)}(E|U,F|U)\otimes_{H^0(U,A)}\Hom_{\Dpx(U)}(F|U,E|U)
  \longrightarrow H^0(U,A),
\]
que nous appellerons \emph{cup-produit}\footnote{même remarque que pour la définition de \(\Tr_{E|U}\)}.

Le lien entre trace et cup-produit s'exprime dans la

\textbf{Proposition 8.3.} \emph{Dans la situation de} (8.2), \emph{pour} \(u\in\Hom(E|U,F|U)\) \emph{et} \(v\in\Hom(F|U,E|U)\), \emph{on a}
\[
  \Tr_{E|U}(vu)=\Tr_{F|U}(uv)=\langle u,v\rangle.
\]
\emph{En particulier}, \emph{si} \(E=F\), \emph{on a}
\[
  \Tr_{E|U}(u)=\langle u,\Id\rangle=\langle\Id,u\rangle.
\]

\emph{Preuve.} Pour simplifier, nous supposerons \(U=S\). On a donc deux flèches \(u:A\longrightarrow E^\vee\otL F\) et \(v:A\longrightarrow F^\vee\otL E\), d'où
\[
  u\otL v:A\longrightarrow E^\vee\otL F\otL F^\vee\otL E.
\]
D'autre part, on a \(vu:A\to E^\vee\otL E\) et \(uv:A\to F^\vee\otL F\). La proposition résulte de la commutativité des diagrammes suivants (laquelle est de vérification «~purement mécanique~») :
\[
\begin{tikzcd}[column sep=large,row sep=large]
 A \arrow[r,"u\otL v"] \arrow[dr,"vu"'] & E^\vee\otL F\otL F^\vee\otL E \arrow[d,"\Id\otL\tau_F"] \\
 & E^\vee\otL E
\end{tikzcd}
\qquad
\begin{tikzcd}[column sep=large,row sep=large]
 A \arrow[r,"u\otL v"] \arrow[dr,"uv"'] & E^\vee\otL F\otL F^\vee\otL E \arrow[d,"\tau_E\otL\Id"] \\
 & F^\vee\otL F
\end{tikzcd}
\]
\[
\begin{tikzcd}[column sep=
%%%%%%%%%%  scan idx 170  |  folio 157  %%%%%%%%%%
large,row sep=large]
E^\vee\otL F\otL F^\vee\otL E \arrow[r,"\Id\otL\tau_E"] \arrow[d,"\Id\otimes\tau_F"'] & F^\vee\otL F \arrow[d,"\tau_F"]\\
E^\vee\otL E \arrow[r,"\tau_E"'] & A .
\end{tikzcd}
\]

\textbf{Corollaire 8.4.} \(E,F,U\) \emph{étant donnés comme en} (8.2), \emph{soient} \(u\in\Hom(E|U,E|U)\), \(s:E|U\to F|U\) \emph{un isomorphisme}, \emph{d'où} \(sus^{-1}\in\Hom(F|U,F|U)\). \emph{On a alors}
\[
  \Tr_{F|U}(sus^{-1})=\Tr_{E|U}(u).
\]

\emph{Preuve.} C'est immédiat.

\textbf{Remarque 8.5.} Si \(\mathcal C\) est une \(S\)-catégorie fibrée, notons \(E(\mathcal C)\) la \(S\)-catégorie fibrée dont la fibre en \(U\in\Ob S\) est la catégorie ayant pour objets les couples \((L,u)\) où \(L\in\Ob\mathcal C_U\) et \(u\in\operatorname{End}_S(L)\) et pour flèches de \((L,u)\) à \((M,v)\) les flèches \(f:L\to M\) telles que \(fu=vf\). On peut, d'après (8.4), considérer l'homomorphisme trace comme un \(S\)-foncteur de \(E(\Dpx(S)_{\underline{\parf}})^{\mathrm{is}}\) dans \(\underline A\) (où l'exposant \(\mathrm{is}\) signifie que les seules flèches considérées sont les isomorphismes, et \(\underline A\) est la \(S\)-catégorie à fibres discrètes définie par \(A\) (cf. (6.6.a))). Comme tout complexe parfait est localement de tor-dimension finie, le \(S\)-foncteur d'oubli
\[
  E(\Dpx(S)_{\underline{\parf}})^{\mathrm{is}}
  \longrightarrow
  E(\Dpx(S)_{\parf})^{\mathrm{is}}
\]
induit un isomorphisme sur les champs associés. Par suite, l'homomorphisme
%%%%%%%%%%  scan idx 171  |  folio 158  %%%%%%%%%%
trace se prolonge de manière unique en un \(S\)-foncteur de \(E(\Dpx(S)_{\parf})^{\mathrm{is}}\) dans \(\underline A\).

Pour \(E\in\Ob\Parf(U)\) (\(U\in\Ob S\)), on note encore
\[
  \Tr_E:\Hom(E,E)\longrightarrow H^0(U,A)
\]
l'homomorphisme correspondant, qui est un homomorphisme de \(H^0(U,A)\)-modules. On voit de même que la définition du cup-produit (8.2), ainsi que les assertions (8.3) et (8.4) se généralisent au cas de complexes parfaits (non nécessairement de tor-dimension finie).

\textbf{Proposition 8.6.} \emph{Soient} \(U\in\Ob S\), \(E,F\in\Parf(U)\), \(u\in\Hom(E,E)\), \(v\in\Hom(F,F)\), \emph{d'où} \(u\otL v\in\Hom(E\otL F,E\otL F)\). \emph{On a}
\[
  \Tr_{E\otL F}(u\otL v)=\Tr_E(u)\,\Tr_F(v).
\]

\emph{Preuve.} La question étant locale, on peut supposer grâce à (8.4) et (4.10), que \(u\) et \(v\) sont des flèches de \(K^b(\mathcal C_0)\). Utilisant (8.1.2), on se ramène au cas où \(u\) et \(v\) sont des flèches de \(\mathcal C_0\). On termine grâce à (8.1.1).

\section*{BIBLIOGRAPHIE}
\begin{description}[leftmargin=2.5cm,style=nextline]
\item[{[H]}]
%%%%%%%%%%  scan idx 172  |  folio 159  %%%%%%%%%%
Hartshorne, R. Residues and Duality, Lecture Notes n\textsuperscript{o}~20, Springer, 1966.
\item[{[V]}] Verdier, J-L. Catégories dérivées, I.H.E.S., 1963.
\item[{[G]}] Giraud, J. Cohomologie non abélienne, à paraître dans North Holl. Pub. Comp.
\item[{SGA 4}] Cohomologie étale des schémas, par M. Artin, A. Grothendieck, et J-L. Verdier, à paraître dans North Holl. Pub. Comp.
\item[{SGA 5}] Cohomologie \(\ell\)-adique et fonctions L, par A. Grothendieck, à paraître dans North Holl. Pub. Comp.
\item[{[1]}] Hakim, M. Topos annelés et schémas relatifs, thèse (en préparation).
\item[{[2]}] Verdier, J-L. Les catégories dérivées des catégories abéliennes, à paraître dans North Holl. Pub. Comp.
\item[{[MULT]}] Cartan, H. et Eilenberg, S. Homological Algebra, Princeton, 1956.
\end{description}



\providecommand{\cS}{\mathcal S}
\providecommand{\cCz}{\mathcal C_0}
\providecommand{\cCone}{\mathcal C_1}
\providecommand{\cA}{\mathcal A}
\providecommand{\cB}{\mathcal B}
\providecommand{\cL}{\mathcal L}
\providecommand{\cI}{\mathcal I}
\providecommand{\cJ}{\mathcal J}
\providecommand{\Db}{\mathrm D}
\providecommand{\Kb}{\mathrm K}
\providecommand{\Dminus}{\mathrm D^{-}}
\providecommand{\Kminus}{\mathrm K^{-}}
\providecommand{\Kbd}{\mathrm K^{b}}
\providecommand{\Mod}{\operatorname{Mod}}
\providecommand{\Lib}{\operatorname{Lib}}
\providecommand{\Loclib}{\operatorname{Loclib}}
\providecommand{\Qcoh}{\operatorname{Qcoh}}
\providecommand{\Coh}{\operatorname{Coh}}
\providecommand{\Ext}{\operatorname{Ext}}
\providecommand{\Hom}{\operatorname{Hom}}
\renewcommand{\Ob}{\operatorname{ob}}
\providecommand{\im}{\operatorname{Im}}
\providecommand{\coker}{\operatorname{Coker}}
\providecommand{\coh}{\operatorname{coh}}
\providecommand{\parf}{\operatorname{parf}}
\providecommand{\amp}{\operatorname{amp}}
\providecommand{\RHom}{\operatorname{RHom}}


%%%%%%%%%%  scan idx 173  |  folio 160  %%%%%%%%%%
\section*{\underline{Exposé II}\\[3pt] \underline{\underline{EXISTENCE DE RÉSOLUTIONS GLOBALES}}}
\begin{center}
\underline{par L. ILLUSIE}
\end{center}

\section*{1. Critères généraux de globalisation}

\textbf{1.0.} Dans ce numéro, on se donne un site $S$, une $S$-catégorie abélienne plate $\mathcal C$ (I 1.1), et une sous-$S$-catégorie additive strictement pleine $\mathcal C_0$ (I 1.1). On suppose $\mathcal C_0$ localement quasi-relevable dans $\mathcal C$ (I 1.2) et localement quasi-stable par noyau d'épimorphisme (I 1.2).

\textbf{Proposition 1.1.} Soit $S\in\Ob S$, et supposons que :
\begin{enumerate}[label=(\roman*)]
\item $\mathcal C_{0S}$ est quasi-relevable dans $\mathcal C_S$ (I 1.2) ($\mathcal C_S$ étant considérée comme catégorie fibrée sur un site ponctuel) ;
\item tout objet de $\mathcal C_S$ qui est de $\mathcal C_0$-type fini (I 1.2) (resp. de $\mathcal C_0$-présentation finie (I 2.8)) est de $\mathcal C_{0S}$-type fini.
\end{enumerate}
Alors :
\begin{enumerate}[label=\alph*)]
\item $\mathcal C_{0S}$ est quasi-stable par noyau d'épimorphisme (I 1.2) ($\mathcal C_{0S}$ étant considérée comme catégorie fibrée sur un site ponctuel).
\item Soient $E\in\Ob \mathrm D^{-}(\mathcal C_S)$ et $n\in\mathbb Z$. Si $E$ est $n$-$\mathcal C_0$-pseudo-cohérent (I 2.3), $E$ est $n$-$\mathcal C_{0S}$-pseudo-cohérent (resp. $(n+1)$-$\mathcal C_{0S}$-pseudo-cohérent). En d'autres termes, on a
\[
\mathrm D^{-}(\mathcal C_S)_{n-\mathcal C_0\text{-coh}}
=
\mathrm D^{-}(\mathcal C_S)_{n-\mathcal C_{0S}\text{-coh}}
\]
(resp.
\[
\mathrm D^{-}(\mathcal C_S)_{n-\mathcal C_0\text{-coh}}
\subset
\mathrm D^{-}(\mathcal C_S)_{(n+1)-\mathcal C_{0S}\text{-coh}}
\subset
\mathrm D^{-}(\mathcal C_S)_{(n+1)-\mathcal C_0\text{-coh}},
\]
inclusions de sous-catégories strictement pleines de $\mathrm D(\mathcal C_S)$, l'inclusion de droite étant d'ailleurs triviale).
\item Pour que $E\in\Ob\mathrm D^{-}(\mathcal C_S)$ soit pseudo-cohérent relativement à $\mathcal C_0$ (I 2.3), il faut et il suffit qu'il le soit relativement à $\mathcal C_{0S}$ ; autrement dit, on a
\[
\mathrm D^{-}(\mathcal C_S)_{\mathcal C_0\text{-coh}}
=
\mathrm D^{-}(\mathcal C_S)_{\mathcal C_{0S}\text{-coh}}.
\]
En outre, le foncteur canonique
\[
\mathrm K^{-}(\mathcal C_{0S})/\mathrm K^{-,\varnothing}(\mathcal C_{0S})
\longrightarrow \mathrm D(\mathcal C_S),
\]
où $\mathrm K^{-,\varnothing}(\mathcal C_{0S})$ désigne la sous-catégorie épaisse de $\mathrm K^{-}(\mathcal C_{0S})$ formée des

\end{enumerate}
\end{document}
